\subsection{注记与进一步结果}
\label{sec:ch3-notes-further-results}

\subsubsection{参考资料}
\label{subsec:ch3-references}

M. Riesz 定理的证明最早发表于论文 \emph{Sur les fonctions conjuguées} (Math. Zeit. 27 (1927), 218--244), 但更早已在 \emph{Les fonctions conjuguées et les séries de Fourier} (C. R. Acad. Sci. Paris 178 (1924), 1464--1467) 中宣布. 正文中的证明只有在 Marcinkiewicz 插值定理于 1939 年出现后才成为可能. Riesz 的原始证明以及 M. Cotlar 的另一证明, 见下文第~\ref{subsec:ch3-other-riesz-proofs}~节; 后一证明见 \emph{A unified theory of Hilbert transforms and ergodic theorems}, Rev. Mat. Cuyana 1 (1955), 105--167. Zygmund \MBUnresolvedReference{citation}{[22]} 还收录了 A. P. Calderón 的 \emph{On the theorems of M. Riesz and Zygmund}, Proc. Amer. Math. Soc. 1 (1950), 533--535, 以及 P. Stein 的 \emph{On a theorem of M. Riesz}, J. London Math. Soc. 8 (1933), 242--247 中的证明.

Kolmogorov 的定理见 \emph{Sur les fonctions harmoniques conjuguées et les séries de Fourier}, Fund. Math. 7 (1925), 23--28. L. H. Loomis 在 \emph{A note on Hilbert's transform}, Bull. Amer. Math. Soc. 52 (1946), 1082--1086 中给出的另一证明可见 Zygmund \MBUnresolvedReference{citation}{[22]}, 而 L. Carleson 的一个优美证明由 Katznelson \MBUnresolvedReference{citation}{[10]} 给出. 我们的证明使用 Calderón--Zygmund 分解, 其使用方式与\MBUnresolvedReference{chapter}{第 5 章}中处理更一般奇异积分时相同. 它源于 A. P. Calderón 与 A. Zygmund 的证明 \emph{On the existence of certain singular integrals}, Acta Math. 88 (1952), 85--139. 引理~\ref{lem:ch3-cotlar-inequality} 见上引 M. Cotlar 的论文.

\subsubsection{共轭函数与 Fourier 级数}
\label{subsec:ch3-conjugate-function-fourier-series}

本章结果也可对定义在单位圆上的函数建立. 从历史上看, 它们的发展与第~\ref{chap:fourier-series-integrals}~章讨论的 Fourier 级数理论密切相关.

给定 $f\in L^1(\mathbb T)$, 等价地说, 给定 $L^1([0,1])$ 中的函数, 它到单位圆盘的调和延拓由与 Poisson 核
\[
 P_r(t)=\frac{1-r^2}{1-2r\cos(2\pi t)+r^2},
 \qquad 0\le r<1,
\]
作卷积得到.

\hypertarget{chapter-03-page-074}{}
$P_r*f(t)$ 的调和共轭由与共轭 Poisson 核
\[
 Q_r(t)=\frac{2r\sin(2\pi t)}{1-2r\cos(2\pi t)+r^2},
 \qquad 0\le r<1,
\]
作卷积得到. Hilbert 变换的对应物是共轭算子, 定义为
\begin{equation}
\label{eq:ch3-conjugation-operator}
 \widetilde f(x)=\lim_{r\to1^-}Q_r*f(x)
 =\lim_{\varepsilon\to0}\int_{\varepsilon<t<1/2}
 \frac{f(x-t)-f(x+t)}{\tan(\pi t)}\,dt.
\end{equation}
可以证明这两个极限对几乎所有 $x$ 都存在.

M. Riesz 与 Kolmogorov 的定理对共轭算子也成立. 因而通过与第~\ref{sec:ch3-multipliers}~节相似的论证, 可以证明 Fourier 级数的部分和算子 $S_Nf$ 在 $L^p$ 上一致有界, 从而 $S_Nf$ 在 $L^p$ 中收敛到 $f$(参见第~\ref{chap:fourier-series-integrals}~章第~\ref{sec:ch1-norm-convergence}~节).

给定函数 $f\in L^1(\mathbb T)$, 它的共轭 Fourier 级数是三角级数
\[
 \widetilde S[f](x)=\sum_{k=-\infty}^{\infty}
 -i\operatorname{sgn}(k)\widehat f(k)e^{2\pi ikx}.
\]
若 $f$ 的 Fourier 级数写成\eqref{eq:ch1-trigonometric-series} 的形式, 则将正弦项与余弦项的系数互换, 并取其差而非其和, 即得共轭级数:
\[
 \widetilde S[f](x)=\sum_{k=0}^{\infty}
 \bigl(a_k\sin(2\pi kx)-b_k\cos(2\pi kx)\bigr).
\]
当 $\widetilde f\in L^1$ 时, $\widetilde S[f]$ 与 $\widetilde f$ 的 Fourier 级数一致.

即使对连续函数, \eqref{eq:ch3-conjugation-operator} 中的第二个极限之所以存在, 也是因为微妙的抵消, 而不是因为 $t$ 接近零时分子变小. 利用 Baire 纲定理可以证明, 存在连续函数 $f$, 使得对每个 $x\in[0,1]$,
\[
 \int_0^{1/2}\frac{|f(x-t)-f(x+t)|}{\tan(\pi t)}\,dt=\infty.
\]
Hilbert 变换也出现同一现象.

关于这些以及更多结果, 见 Bary \MBUnresolvedReference{citation}{[1]}、Katznelson \MBUnresolvedReference{citation}{[10]}、Koosis \MBUnresolvedReference{citation}{[11]} 和 Zygmund \MBUnresolvedReference{citation}{[21]} 的著作.

\subsubsection{M. Riesz 定理的其他证明}
\label{subsec:ch3-other-riesz-proofs}

M. Riesz 关于共轭函数 $L^p$ 有界性的原始证明以解析函数的幂的解析性为基础. 最容易的情形是 $p=2k$, 其中 $k$ 为整数. 给定单位圆上的函数 $f$, 令 $u=P_r*f$、$v=Q_r*f$; 则 $F(z)=u+iv$ 在单位圆盘内解析, 因而 $F(z)^{2k}$ 也解析. 由 Cauchy 定理,
\[
 \frac1{2\pi i}\int_{|z|=r}\frac{F(z)^{2k}}z\,dz
 =\frac1{2\pi}\int_0^{2\pi}F(re^{it})^{2k}\,dt
 =F(0)^{2k}.
\]

\hypertarget{chapter-03-page-075}{}
若 $F(0)=0$, 对 $F(re^{it})^{2k}$ 取实部, 得
\[
 \int_0^{2\pi}|v(re^{it})|^{2k}\,dt
 \le\sum_{j=1}^k\binom{2k}{2j}
 \int_0^{2\pi}|u(re^{it})|^{2j}|v(re^{it})|^{2k-2j}\,dt.
\]
对右端每一项应用具有适当指数的 Hölder 不等式, 得
\[
 \int_0^{2\pi}|v(re^{it})|^{2k}\,dt
 \le C_k\int_0^{2\pi}|u(re^{it})|^{2k}\,dt,
\]
其中 $C_k$ 与圆的半径无关. 令半径趋于 $1$, $u$ 趋于 $f$, 而 $v$ 趋于 $\widetilde f$. 当 $F(0)=\widehat f(0)\ne0$ 时, 用 $f-\widehat f(0)$ 代替 $f$.

从 $p=2k$ 时的不等式出发, 可以用插值与对偶性得到整个定理. 但 1923 年尚无插值理论可用, 因而 Riesz 随后考虑 $p>1$ 且 $p$ 不是整数的情形. 为使 $(u+iv)^p$ 解析, 他假设 $u>0$(当 $f$ 非负且不恒为零时确实如此), 然后使用了一个巧妙的估计. 由于技术原因, 该论证对奇整数 $p$ 不适用, 这些 $p$ 值的不等式则由对偶性得到. 在一封写于 1923 年 11 月的信中, Riesz 向 Hardy 解释了自己如何得到这一证明, 并向他展示了关键细节; 该信转载于 M. L. Cartwright, \emph{Manuscripts of Hardy, Littlewood, M. Riesz and Titchmarsh}, Bull. London Math. Soc. 14 (1982), 472--532.

由 $(u+iv)^2=u^2-v^2+2iuv$ 解析这一事实, 得
\[
 H\bigl(f^2-(Hf)^2\bigr)=2f\cdot Hf,
\]
从而得到 Cotlar 的公式
\begin{equation}
\label{eq:ch3-cotlar-square-identity}
 (Hf)^2=f^2+2H(f\cdot Hf).
\end{equation}
Cotlar 在其论文中提出了另一种证明方式, 即直接证明
\[
 \widehat{Hf}*\widehat{Hf}
 =\widehat f*\widehat f
 -2i\operatorname{sgn}(\xi)\bigl(\widehat f*\widehat{Hf}\bigr).
\]
由\eqref{eq:ch3-cotlar-square-identity} 可知, 如果 $\|Hf\|_p\le C_p\|f\|_p$, 则
\[
 \|Hf\|_{2p}\le\left(C_p+\sqrt{C_p^2+1}\right)\|f\|_{2p}.
\]
事实上,
\[
\begin{aligned}
 \|Hf\|_{2p}^2
 &=\|(Hf)^2\|_p\\
 &\le\|f^2\|_p+2\|H(f\cdot Hf)\|_p\\
 &\le\|f\|_{2p}^2+2C_p\|f\cdot Hf\|_p\\
 &\le\|f\|_{2p}^2+2C_p\|f\|_{2p}\|Hf\|_{2p},
\end{aligned}
\]

\hypertarget{chapter-03-page-076}{}
解这个二次不等式即得所需结论.

从\eqref{eq:ch3-hilbert-l2-isometry} 出发反复应用这一不等式, 例如可得对整数 $k\ge1$,
\[
 \|Hf\|_{2^k}\le(2^k-1)\|f\|_{2^k}.
\]
再由 Riesz--Thorin 插值定理(定理~\ref{thm:ch1-riesz-thorin}), 得
\[
 \|Hf\|_p\le C_p\|f\|_p,
 \qquad p\ge2.
\]
像定理~\ref{thm:ch3-riesz-kolmogorov} 的证明那样使用对偶性, 可得到 $p<2$ 时的 Riesz 定理.

$p=2$ 时还有下面这个不使用 Fourier 变换的初等实变量证明:
\[
 \|H_\varepsilon f\|_2^2
 =\langle f,-H_\varepsilon H_\varepsilon f\rangle
 \le\|f\|_2\|H_\varepsilon H_\varepsilon f\|_2.
\]
$H_\varepsilon H_\varepsilon$ 的核是 $y^{-1}\chi_{\{|y|>\varepsilon\}}$ 与自身的卷积, 它可以显式算出并证明属于 $L^1$. 通过伸缩论证, 只需考虑 $\varepsilon=1$ 的情形. E. M. Dyn'kin 在 \emph{Methods of the theory of singular integrals: Hilbert transform and Calderón--Zygmund theory}, \emph{Commutative Harmonic Analysis}, I, V. Havin and N. Nikolskii 编, pp. 167--259, Springer-Verlag, Berlin, 1991 中将这个证明归于 N. Lusin (1951). Schur 更早曾对算子的离散版本使用类似论证.

另一个在 $p=2$ 时使用 Cotlar 引理的实变量证明见\MBUnresolvedReference{section}{第 9 章第 1 节}.

\subsubsection{常数的大小}
\label{subsec:ch3-size-of-constants}

Kolmogorov 与 M. Riesz 定理中的最佳常数是已知的. 对强 $(p,p)$ 不等式, $1<p<\infty$, 最佳常数为
\[
 C_p=\tan\frac{\pi}{2p},\quad1<p\le2;
 \qquad
 C_p=\cot\frac{\pi}{2p},\quad2\le p<\infty.
\]
S. K. Pichorides 在 \emph{On the best values of the constants in the theorems of M. Riesz, Zygmund and Kolmogorov}, Studia Math. 44 (1972), 165--179 中首先证明这一结果; L. Grafakos 后来在 \emph{Best bounds for the Hilbert transform on $L^p(\mathbb R^1)$}, Math. Res. Let. 4 (1997), 469--471 中给出了简单得多的证明.

在上面 Cotlar 给出的 M. Riesz 定理证明中, 得到估计
\[
 C_{2p}\le C_p+\sqrt{C_p^2+1}.
\]
若从 $C_2=1$ 出发, 所得到的 $C_{2^k}$ 上界是锐的. I. Gohberg 与 N. Krupnik 在 \emph{Norm of the Hilbert transformation in the space $L^p$}, Funct. Anal. Appl. 2 (1968), 180--181 中首先注意到这一事实.

\hypertarget{chapter-03-page-077}{}
弱 $(1,1)$ 不等式中的最佳常数具有更复杂的表达式:
\[
 C_1=\frac{1+\dfrac1{3^2}+\dfrac1{5^2}+\cdots}
 {1-\dfrac1{3^2}+\dfrac1{5^2}-\cdots}.
\]
B. Davis 使用 Brownian 运动首先证明了这一结果, 见 \emph{On the weak type $(1,1)$ inequality for conjugate functions}, Proc. Amer. Math. Soc. 44 (1974), 307--311. A. Baernstein 后来在 \emph{Some sharp inequalities for conjugate functions}, Indiana Univ. Math. J. 27 (1978), 833--852 中给出了一个非概率证明.

\subsubsection{$L^1$ 上的 Hilbert 变换}
\label{subsec:ch3-hilbert-l1}

正如第~\ref{sec:ch3-riesz-kolmogorov-theorems}~节指出的, 若 $f\in L^1$, 则 $Hf$ 未必属于 $L^1$. 但是, 与极大函数的情形不同, 存在 $f\in L^1$ 使得 $Hf\in L^1$. 例如, 若 $f=\chi_{(0,1)}-\chi_{(-1,0)}$, 则
\[
 Hf(x)=\log\frac{|x^2-1|}{x^2}
\]
是可积的.

当 $f\in L^1$ 时, $Hf$ 可积的一个简单必要条件是 $\widehat f(0)=\int f=0$. 若 $f$ 满足恒等式 $\widehat{Hf}(\xi)=-i\operatorname{sgn}(\xi)\widehat f(\xi)$, 只需注意可积函数的 Fourier 变换连续, 而 $\widehat f(\xi)$ 与 $-i\operatorname{sgn}(\xi)\widehat f(\xi)$ 只有在 $\widehat f(0)=0$ 时才能同时连续. 当 $f\in L^1\cap L^2$ 时, Fourier 变换恒等式成立, 因而这个论证显然适用. 但可以证明, 当 $f,Hf\in L^1$ 时该恒等式也成立. 为此, 若能对每个 $\varphi\in\mathcal S$ 证明
\[
 \varphi*(Hf)=H(\varphi*f),
\]
则对两边取 Fourier 变换并化简即可. A. P. Calderón 与 O. Capri 在 \emph{On the convergence in $L^1$ of singular integrals}, Studia Math. 78 (1984), 321--327 中按照处理更一般奇异积分的方法证明了这一等式.

同样, 当 $f,Hf\in L^1$ 时, 有恒等式 $H(Hf)=-f$. 我们在\eqref{eq:ch3-hilbert-square-minus-identity} 中首先对 $L^2$ 函数给出这个恒等式, 利用本章结果可将它延拓到 $L^p$, $1<p<\infty$. 当 $f$ 与 $Hf$ 可积时, E. Hille 与 J. D. Tamarkin 在 \emph{On the absolute integrability of the Fourier transform}, Fund. Math. 25 (1935), 329--352 中给出了使用复分析的证明. J. F. Toland 的论文 \emph{A few remarks about the Hilbert transform}, J. Funct. Anal. 145 (1997), 151--174 考察了这一问题, 并给出了更多历史资料.

一般而言 $Hf\notin L^1$, 因而甚至不能定义 $Hf$ 的 Hilbert 变换. 不过, 如果用一种更一般的积分取代 Lebesgue 积分, 就可以定义 $H(Hf)$ 并得到上述恒等式(见 Zygmund \MBUnresolvedReference{citation}{[21, Chapter 7]}). 上引 Toland 的论文使用了这种推广的积分概念.

\hypertarget{chapter-03-page-078}{}
\subsubsection{$L\log L$ 估计}
\label{subsec:ch3-l-log-l-estimates}

下述结果刻画 Hilbert 变换的局部可积性; 它与极大函数的定理~\ref{thm:ch2-l-log-l-characterization} 类似.

\begin{Theorem}
\label{thm:ch3-l-log-l-characterization}
设 $B$ 与 $C$ 是满足 $\overline B\subset C$ 的两个开球.

\begin{enumerate}
\item 若 $f\log^+|f|\in L^1(C)$, 则 $Hf\in L^1(B)$.
\item 反之, 若 $f\ge0$ 且 $Hf\in L^1(C)$, 则 $f\log^+|f|\in L^1(B)$.
\end{enumerate}
\end{Theorem}

定理第一部分归功于 A. P. Calderón 与 A. Zygmund(见上引论文). 利用 Orlicz 空间可以加强这一结论(参见第~\ref{subsec:ch2-l-log-l}~节以及 Zygmund \MBUnresolvedReference{citation}{[21, Chapter 4]}). 第二部分归功于 E. M. Stein 的 \emph{Note on the class $L\log L$}, Studia Math. 32 (1969), 305--310. 假设 $f\ge0$ 可以替换为一个更弱的条件, 它在某种意义下度量 $P_t*f$ 与一个正调和函数的差异程度. 见 J. Brossard 与 L. Chevalier 的论文 \emph{Classe $L\log L$ et densité de l'intégrale d'aire dans $\mathbb R_+^{n+1}$}, Ann. of Math. 128 (1988), 603--618.

\subsubsection{平移不变算子与乘子}
\label{subsec:ch3-translation-invariant-operators}

令 $\tau_h$ 为平移算子: $\tau_hf(x)=f(x-h)$. 如果算子 $T$ 与 $\tau_h$ 交换, 即
\[
 \tau_h\circ T=T\circ\tau_h,
 \qquad h\in\mathbb R^n,
\]
就称 $T$ 平移不变. 例如 Hilbert 变换是平移不变的. 这类算子由下述结果完全刻画.

\begin{Theorem}
\label{thm:ch3-translation-invariant-convolution}
设线性算子 $T$ 平移不变, 并且对某对 $(p,q)$, $1\le p,q\le\infty$, 从 $L^p(\mathbb R^n)$ 有界映入 $L^q(\mathbb R^n)$. 则存在唯一的缓增分布 $K$, 使得
\[
 Tf=K*f,
 \qquad f\in\mathcal S.
\]
\end{Theorem}

若这样的 $T$ 从 $L^p$ 有界映入 $L^q$ 且不为零, 则 $p$ 不能大于 $q$. L. Hörmander 在 \emph{Estimates for translation invariant operators in $L^p$-spaces}, Acta Math. 104 (1960), 93--139 中给出了这一事实的一个优美而极其简单的证明. 首先注意到
\[
 \lim_{|h|\to\infty}\|f+\tau_hf\|_p=2^{1/p}\|f\|_p.
\]
以 $\|T\|_{p,q}$ 表示 $T$ 从 $L^p$ 到 $L^q$ 的算子范数. 因为 $T$ 与平移交换,
\[
 \|Tf+\tau_hTf\|_q\le\|T\|_{p,q}\|f+\tau_hf\|_p.
\]
令 $h$ 趋于无穷, 得
\[
 \|Tf\|_q\le2^{1/p-1/q}\|T\|_{p,q}\|f\|_p,
\]
而这只有在 $p\le q$ 时才可能成立.

显然, 若 $Tf=K*f$, 则 $T$ 线性且平移不变. 由 Fourier 变换的性质,
\[
 \widehat{Tf}=\widehat K\,\widehat f,
\]
这说明至少当 $\widehat K\in L^\infty$ 时, $\widehat K$ 是第~\ref{sec:ch3-multipliers}~节定义意义下 $T$ 的乘子.

\hypertarget{chapter-03-page-079}{}
但是, 若 $T$ 在 $L^p$ 上有界, 则由对偶性, 它的伴随算子是与核 $\widetilde K(x)=\overline{K(-x)}$ 的卷积, 因而在 $L^{p'}$ 上有界. 所以由插值可知 $T$ 在 $L^2$ 上有界, 条件 $\widehat K\in L^\infty$ 是必要的, 从而我们关于乘子 $m\in L^\infty$ 的假设是合理的.

刻画定义有界算子的分布 $K$(等价地说, 刻画 $L^p$ 乘子)是一个有趣的问题, 但完整答案只在两种情形下已知.

\begin{Proposition}
\label{prop:ch3-l1-l2-convolution-characterization}
下列结论成立:

\begin{enumerate}
\item $T$ 在 $L^2(\mathbb R^n)$ 上有界, 当且仅当 $\widehat K\in L^\infty$; $T$ 在 $L^2$ 上的算子范数等于 $\widehat K$ 的 $L^\infty$ 范数.
\item $T$ 在 $L^1(\mathbb R^n)$ 上有界, 当且仅当 $K$ 是有限 Borel 测度; $T$ 的算子范数等于该测度的全变差.
\end{enumerate}
\end{Proposition}

注意, 因为 $\operatorname{p.v.}\,1/x$ 不是测度, 命题~\ref{prop:ch3-l1-l2-convolution-characterization} 给出了 Hilbert 变换在 $L^1$ 上不有界的另一个证明.

关于上述所有结果, 见 Stein 与 Weiss \MBUnresolvedReference{citation}{[18, Chapter 1]}.

\subsubsection{Fourier 级数的乘子}
\label{subsec:ch3-fourier-series-multipliers}

三角级数的乘子恰好是有界序列 $\{\lambda(k)\}_{k\in\mathbb Z}$. 对 $f\in L^2(\mathbb T)$, 相伴算子定义为
\[
 T_\lambda f(x)=\sum_{k\in\mathbb Z}\lambda(k)\widehat f(k)e^{2\pi ikx}.
\]
对多重 Fourier 级数, 除了取 $k\in\mathbb Z^n$ 外, 定义相同. $T_\lambda$ 可写成 $f$ 与一个 Fourier 系数为 $\{\lambda(k)\}$ 的分布的卷积, 并且有与命题~\ref{prop:ch3-l1-l2-convolution-characterization} 类似的结果.

如果知道 $L^p(\mathbb R)$ 上一个乘子的有界性, 就可以把乘子限制到整数上, 从而推导级数乘子的结果.

\begin{Theorem}
\label{thm:ch3-multiplier-restriction-to-integers}
设 $1\le p\le\infty$, 算子 $Tf=K*f$ 在 $L^p(\mathbb R)$ 上有界. 若 $\widehat K$ 在 $\mathbb Z$ 的每一点连续, 且 $\lambda(k)=\widehat K(k)$, 则与这个序列相伴的算子 $T_\lambda$ 在 $L^p(\mathbb T)$ 上有界, 并且
\[
 \|T_\lambda\|\le\|T\|.
\]
\end{Theorem}

还有一个类似的 $n$ 维结果. 证明见 Stein 与 Weiss \MBUnresolvedReference{citation}{[18, Chapter 6]}.

从定理~\ref{thm:ch3-multiplier-restriction-to-integers} 出发, 可以由第~\ref{sec:ch3-multipliers}~节中 $S_R$ 的相应结果推导 Fourier 级数部分和算子 $S_N$ 的结果,

\hypertarget{chapter-03-page-080}{}
特别是可得到 Fourier 级数部分和依 $L^p$ 范数收敛(参见第~\ref{chap:fourier-series-integrals}~章第~\ref{sec:ch1-norm-convergence}~节).

\subsubsection{特征函数的 Hilbert 变换}
\label{subsec:ch3-characteristic-function-hilbert-transform}

若 $A$ 是 $\mathbb R$ 中的有限测度子集, 则可以证明
\[
 \bigl|\{x\in\mathbb R:|H(\chi_A)(x)|>\lambda\}\bigr|
 =\frac{2|A|}{\sinh(\pi\lambda)}.
\]
E. M. Stein 与 G. Weiss 在 \emph{An extension of a theorem of Marcinkiewicz and some of its applications}, J. Math. Mech. 8 (1959), 263--284 中首先证明这一结果. Zygmund \MBUnresolvedReference{citation}{[22, p. 15]} 给出了一个简单证明.
